\section*{Chapitre \ref{ch:b2:basins-compared} --- Bassins sédimentaires comparés}

\begin{solution}{exo:b2:basins-compared:1}
(a) Intracratonique : enfoncement lent, par pulsations, pendant des centaines de millions d'années, sans faille
majeure~; sa cause est débattue. (b) Rift : amincissement mécanique de la croûte pendant l'étirement, puis
contraction thermique à mesure que la lithosphère se refroidit. (c) Marge passive : subsidence thermique héritée du stade de rift,
amplifiée par la charge d'un épais coin de sédiments. (d) et (e) Avant-pays : la lithosphère fléchie vers le bas par la charge
de la chaîne de chevauchement, des Alpes et de l'Himalaya.
\end{solution}

\begin{solution}{exo:b2:basins-compared:2}
(a) $Y = 3000\times(3330 - 2300)/(3330 - 1030) = 3000\times1030/2300 = \qty{1343}{m}$.
(b) $1343 + 200 = \qty{1543}{m}$. (c) Sans sa charge, le socle se trouverait à \qty{1343}{m}~:
$1 - 1343/3000 = \qty{55}{\percent}$ de l'épaisseur est de la place créée par la charge.
\end{solution}

\begin{solution}{exo:b2:basins-compared:3}
(a) $\beta = 35/25 = 1.4$. (b) $125/1.4 = \qty{89.3}{km}$. (c) $S_i = 4334\times(1 - 1/1.4) =
\qty{1238}{m}$.
\end{solution}

\begin{solution}{exo:b2:basins-compared:4}
(a) $\tau = (1.25\times10^{5})^2/(\pi^2\times8\times10^{-7}) = \qty{1.979e15}{s}$, soit
$1.979\times10^{15}/3.156\times10^{13} = \qty{62.7}{Ma}$. (b) $62.7\times0.8 = \qty{50.2}{Ma}$.
(c) $62.7\times(100/125)^2 = \qty{40.1}{Ma}$. (d) $1 - e^{-1} = \qty{63}{\percent}$~; $1 -
e^{-3} = \qty{95}{\percent}$.
\end{solution}

\begin{solution}{exo:b2:basins-compared:5}
(a) $0.63\,e^{-1000/1960} = 0.378$; $0.63\,e^{-2000/1960} = 0.227$. (b) $T_s = 1000 -
0.63\times1960\times(0.6004 - 0.3605) = \qty{704}{m}$. (c) $m = 2700\times704 + 1030\times296 =
\qty{2.206e6}{kg.m^{-2}}$; $\bar\rho = \qty{2206}{kg.m^{-3}}$.
\end{solution}

\begin{solution}{exo:b2:basins-compared:6}
(a) $z_2' \leftarrow 704 + 1234.8\,(1 - e^{-z_2'/1960})$~: 1077, 1226, 1278, 1296, 1301, 1303,
1304~; la couche est épaisse de \qty{1304}{m}. (b) $(2700\times704 + 1030\times600)/1304 =
\qty{1932}{kg.m^{-3}}$. (c) De \qty{304}{m}, soit \qty{30}{\percent}~: l'eau que la compaction
a chassée.
\end{solution}

\begin{solution}{exo:b2:basins-compared:7}
(a) $1 - e^{-100/62.7} = 0.7971$. Pour $\beta = 1.5$~: $3207\times0.4135\times0.7971 =
\qty{1057}{m}$~; pour $\beta = 2$~: $3207\times0.6366\times0.7971 = \qty{1627}{m}$. (b) Oui~: l'intervalle 1100 à
\qty{1400}{m} se situe entre les deux. (c) $\frac{\beta}{\pi}\sin\frac{\pi}{\beta} = 1250/(3207\times0.7971)
= 0.4890$, donc $\beta = 1.6 + 0.1\times(0.4890 - 0.4705)/(0.5204 - 0.4705) = 1.64$, dans la
fourchette de 50 à \qty{100}{\percent} d'étirement déduite pour le bassin.
\end{solution}

\begin{solution}{exo:b2:basins-compared:8}
(a) La masse volumique du remplissage remplace $\rho_w$ au dénominateur. (b) Le numérateur vaut
$125\,000\times(1/3)\times74.69 = \num{3.112e6}$~; $S_i = 3.112\times10^{6}/(3184.4 - 2300) =
\qty{3519}{m}$. (c) Sous l'eau, $3.112\times10^{6}/2154.4 = \qty{1445}{m}$~: le bassin rempli de sédiment est
2.44 fois plus profond.
\end{solution}

\begin{solution}{exo:b2:basins-compared:9}
(a) $2171 - 100\times3330/2300 = 2171 - 145 = \qty{2026}{m}$. (b) $Y$ varie mètre pour mètre
avec $W_d$~: de 1971 à \qty{2371}{m}. (c) La profondeur d'eau, $\pm\qty{200}{m}$, contre
\qty{145}{m} pour le niveau de la mer.
\end{solution}

\begin{solution}{exo:b2:basins-compared:10}
(a) $S_i$ a le signe du crochet de son numérateur, puisque le dénominateur, $3184.4 - 1030$, est
positif~; il est négatif, un soulèvement, lorsque le crochet l'est. (b) $(\rho_m - \rho_c)t_c/a =
\alpha T_1\rho_m/2$ donne $t_c = a\,\alpha T_1\rho_m/(2(\rho_m - \rho_c)) = 125\times0.04372\times
3330/1060 = \qty{17.17}{km}$. (c) $t_c = 17.17/(1 - 0.04372\times17.17/250) = \qty{17.22}{km}$.
(d) Là où la croûte est plus mince qu'environ \qty{17}{km} avant l'étirement : dans la lithosphère océanique,
ou dans une croûte continentale déjà amincie par un rifting antérieur.
\end{solution}

\begin{solution}{exo:b2:basins-compared:11}
(a) Avant l'étirement, la colonne contient la lithosphère de masse $M = \rho_c t_c(1 - \alpha T_1
t_c/(2a)) + \rho_m(a - t_c)(1 - \alpha T_1(a + t_c)/(2a))$~; après, $\rho_w S_i + M/\beta +
\rho_m(1 - \alpha T_1)(a - S_i - a/\beta)$. En égalant les deux, $S_i[\rho_m(1 - \alpha T_1) - \rho_w] =
(1 - 1/\beta)[\rho_m(1 - \alpha T_1)a - M]$, et en développant $M$, $\rho_m(1 - \alpha T_1)a - M =
(\rho_m - \rho_c)t_c(1 - \alpha T_1 t_c/(2a)) - \rho_m\alpha T_1 a/2$~: la formule. (b) Le profil final est $T_1 z/a$~; le profil initial le dépasse
de $T_1 z(\beta - 1)/a$ au-dessus de $a/\beta$ et de $T_1(1 - z/a)$ en dessous. Les intégrales valent $T_1
a(\beta - 1)/(2\beta^2)$ et $T_1 a(\beta - 1)^2/(2\beta^2)$, dont la somme est $T_1 a(1 - 1/\beta)/2$.
La contraction $\alpha$ fois cette quantité, chargée par l'eau, donne l'expression. (c) $3.28\times10^{-5}
\times3330\times1333\times125\,000\times0.5/(2\times2300) = \qty{1978}{m}$, contre $3207\times
2/\pi = \qty{2042}{m}$~: le premier terme surestime la subsidence finale de \qty{64}{m},
\qty{3}{\percent}.
\end{solution}

\begin{solution}{exo:b2:basins-compared:12}
(a) $3207\times0.6366\times(1 - e^{-150/62.7}) = 2042\times0.9086 = \qty{1855}{m}$. (b) $1855\times
2300/(3330 - 2400) = \qty{4588}{m}$. (c) De 6.9 à \qty{13.8}{km}. (d) Le modèle néglige plusieurs
facteurs susceptibles d'approfondir le bassin : un manteau lithosphérique plus étiré que la croûte, une réponse
régionale (flexurale) à la charge du sédiment, un soulèvement transitoire de la marge pendant le rifting
suivi d'un enfoncement, ou un apport de sédiment si grand que le bassin fut comblé puis chargé
au-delà du niveau de la mer sur la plate-forme.
\end{solution}

\begin{solution}{pb:b2:basins-compared:1}
\textbf{1.} $0.49\,e^{-4420/3704} = 0.149$ au sommet~; $0.49\,e^{-6030/3704} = 0.096$ à la base.
\textbf{2.} $T_s = 1610 - 0.49\times3704\times(0.3032 - 0.1963) = \qty{1416}{m}$.
\textbf{3.} Argilite : $1110 - 0.63\times1960\times(e^{-3310/1960} - e^{-4420/1960}) = 1110 - 1234.8
\times(0.1847 - 0.1049) = \qty{1011}{m}$.
\textbf{4.} Sable argileux : $2680\times1004 + 1030\times696 = \qty{3.408e6}{kg.m^{-2}}$~; craie :
$2710\times1409 + 1030\times201 = \qty{4.025e6}{kg.m^{-2}}$~; argilite : $2700\times1011 + 1030\times99 =
\qty{2.832e6}{kg.m^{-2}}$~; sable : $2650\times1416 + 1030\times194 = \qty{3.952e6}{kg.m^{-2}}$. Total
\qty{1.4217e7}{kg.m^{-2}}, $\bar\rho_s = 14\,216\,910/6030 = \qty{2358}{kg.m^{-3}}$.
\textbf{5.} $z_2' \leftarrow 1416 + 1815\,(1 - e^{-z_2'/3704})$~: 1993, 2171, 2221, 2235, 2238, 2239,
2239~: \qty{2239}{m}.
\textbf{6.} $(2650\times1416 + 1030\times823)/2239 = \qty{2055}{kg.m^{-3}}$.
\textbf{7.} $Y = 2239\times(3330 - 2055)/2300 + 200 = 1241 + 200 = \qty{1441}{m}$.
\textbf{8.} $Y = 6030\times(3330 - 2358)/2300 + 100 = 2548 + 100 = \qty{2648}{m}$. Les unités sont
mesurées à leurs profondeurs actuelles : elles sont déjà aussi compactées qu'aujourd'hui.
\textbf{9.} $(2171 - 1441)/50 = \qty{14.6}{m}$ par Ma de 150 à \qty{100}{Ma}~;
$(2453 - 2171)/40 = \qty{7.05}{m}$ de 100 à \qty{60}{Ma}~; $(2648 - 2453)/60 = \qty{3.25}{m}$
depuis~: l'enfoncement ralentit.
\textbf{10.} $6030 + 100 = \qty{6130}{m}$, soit 2.31 fois la subsidence tectonique.
\textbf{11.} $(6130 - 2648)/6130 = \qty{57}{\percent}$.
\textbf{12.} $2171 - 100\times3330/2300 = \qty{2026}{m}$.
\textbf{13.} Le point de fin de rifting, \qty{150}{Ma}~: $S_i = \qty{1441}{m}$, si aucune subsidence thermique
ne s'est produite pendant les \qty{15}{Ma} de rifting.
\textbf{14.} $\beta = 1/(1 - 1441/4334) = 1/0.6675 = 1.50$.
\textbf{15.} $35/1.50 = \qty{23.3}{km}$ de croûte et $125/1.50 = \qty{83.3}{km}$ de lithosphère.
\textbf{16.} $(1.50/\pi)\sin(\pi/1.50) = 0.4775\times0.8660 = 0.4135$~; $3207\times0.4135 =
\qty{1326}{m}$.
\textbf{17.} $S_i = 4334/3 = \qty{1445}{m}$. À \qty{100}{Ma}, $t = \qty{50}{Ma}$~: $1326\times(1 -
e^{-50/62.7}) = \qty{729}{m}$, $Y = \qty{2173}{m}$~; à \qty{60}{Ma}~: $1326\times0.7620 =
\qty{1010}{m}$, $Y = \qty{2455}{m}$~; aujourd'hui~: $1326\times0.9086 = \qty{1205}{m}$, $Y = \qty{2650}{m}$.
\textbf{18.} Le modèle est \qty{2}{m} plus profond en chaque point : 2173 contre 2171, 2455 contre 2453,
2650 contre 2648.
\textbf{19.} Profondeurs d'eau anciennes connues seulement à une centaine de mètres près ou davantage~; niveaux marins passés
différents de l'actuel~; un rifting qui a duré \qty{15}{Ma}, pendant lequel une part de la chaleur était déjà
perdue~; une flexure plutôt qu'une isostasie locale~; un manteau étiré plus ou moins que la croûte~;
une croûte qui n'avait pas \qty{35}{km} d'épaisseur avant le rifting.
\textbf{20.} $1326\times(e^{-150/62.7} - e^{-200/62.7}) = \qty{67}{m}$.
\textbf{21.} $66.6\times2300/(3330 - 2400) = \qty{165}{m}$.
\textbf{22.} $1 - e^{-150/62.7} = \qty{91}{\percent}$.
\textbf{23.} $\tau\ln2 = \qty{43.5}{Ma}$ après la fin du rifting~: à \qty{106.5}{Ma}.
\textbf{24.} $\beta = \textbf{1.50}$~: la croûte a été étirée jusqu'à une fois et demie sa largeur.
\end{solution}
